\documentclass[UTF8,a4paper,11pt]{ctexart}
\usepackage{amsmath,amssymb,bm,geometry,fancyhdr,enumitem}
\geometry{left=22mm,right=22mm,top=22mm,bottom=22mm}
\pagestyle{fancy}
\fancyhf{}
\lhead{现代控制理论·第二周作业标准解答}
\rhead{习题七}
\cfoot{\thepage}
\setlength{\parindent}{2em}
\setlength{\parskip}{0.35em}
\begin{document}
\begin{center}
{\LARGE\bfseries 现代控制理论·第二周作业标准解答}\\[4pt]
{\large 习题七：7-9，7-10(a)(c)，7-12(a)，7-15，7-17}
\end{center}

\section*{7-9}
已知
\[
G(z)=\frac{Y(z)}{U(z)}=\frac{z+1}{z^2-1.4z+0.48},
\]
输入为单位阶跃，故
\[
U(z)=\frac{z}{z-1}.
\]
系统极点满足
\[
z^2-1.4z+0.48=(z-0.8)(z-0.6),
\]
均位于单位圆内，因此可用离散终值定理：
\[
\begin{aligned}
y(\infty)
&=\lim_{z\to1}\frac{z-1}{z}Y(z)\\
&=\lim_{z\to1}\frac{z-1}{z}G(z)\frac{z}{z-1}
=G(1)\\
&=\frac{1+1}{1-1.4+0.48}=25.
\end{aligned}
\]
\[
\boxed{y(\infty)=25}
\]

\section*{7-10(a)}
记连续环节 $F(s)$ 的脉冲传递函数为
\[
F(z)=\mathcal Z\left\{\mathcal L^{-1}[F(s)]\big|_{t=kT}\right\}.
\]
设主采样器输出为 $E(z)$。由前向通道
\[
Y(z)=H(z)E(z).
\]
反馈支路中，$H(s)$ 与 $G_2(s)$ 之间没有采样器，因此二者应先在 $s$ 域相乘后再作 $z$ 变换，记
\[
HG_2(z)=\mathcal Z\{H(s)G_2(s)\},
\]
则
\[
V(z)=HG_2(z)E(z).
\]
经过反馈支路采样器后，$G_1(s)$ 与前级之间被采样器隔开，故
\[
B(z)=G_1(z)V(z)=G_1(z)HG_2(z)E(z).
\]
主采样点满足
\[
E(z)=R(z)-B(z),
\]
故
\[
E(z)\big[1+G_1(z)HG_2(z)\big]=R(z).
\]
于是闭环脉冲传递函数
\[
\boxed{\frac{Y(z)}{R(z)}
=\frac{H(z)}{1+G_1(z)HG_2(z)}}.
\]

\section*{7-10(c)}
定义
\[
G_{01}(z)=\mathcal Z\{H_0(s)G_1(s)\},\qquad
G_{012}(z)=\mathcal Z\{H_0(s)G_1(s)G_2(s)\}.
\]
设数字控制器输入为 $E(z)$，则
\[
D(z)=G_D(z)E(z).
\]
由结构得
\[
C(z)=G_{01}(z)D(z),\qquad
Y(z)=G_{012}(z)D(z).
\]
两个反馈均为负反馈，故
\[
E(z)=X(z)-C(z)-Y(z).
\]
代入可得
\[
E(z)=X(z)-G_D(z)G_{01}(z)E(z)-G_D(z)G_{012}(z)E(z),
\]
即
\[
E(z)=\frac{X(z)}
{1+G_D(z)G_{01}(z)+G_D(z)G_{012}(z)}.
\]
所以
\[
\boxed{
\frac{Y(z)}{X(z)}
=
\frac{G_D(z)G_{012}(z)}
{1+G_D(z)G_{01}(z)+G_D(z)G_{012}(z)}
}.
\]

\section*{7-12(a)}
记
\[
G_1(z)=\mathcal Z\{G_1(s)\},\qquad
G_3(z)=\mathcal Z\{G_3(s)\},\qquad
G_{12}(z)=\mathcal Z\{G_1(s)G_2(s)\}.
\]
由于 $G_1(s)$ 与 $G_2(s)$ 之间没有采样器，内反馈通道应使用整体脉冲传递函数 $G_{12}(z)$。

设第一个采样器输出为 $A(z)$，第二个加法点送入 $G_1(s)$ 的采样序列为 $Q(z)$。内环满足
\[
Q(z)=A(z)-G_{12}(z)Q(z),
\]
故
\[
A(z)=[1+G_{12}(z)]Q(z).
\]
输出采样值为
\[
C(z)=G_1(z)Q(z).
\]
外反馈支路满足
\[
A(z)=R(z)-G_3(z)C(z)
=R(z)-G_3(z)G_1(z)Q(z).
\]
联立得
\[
[1+G_{12}(z)+G_1(z)G_3(z)]Q(z)=R(z).
\]
因此
\[
\boxed{
\frac{C(z)}{R(z)}
=
\frac{G_1(z)}
{1+G_{12}(z)+G_1(z)G_3(z)}
}.
\]

\section*{7-15}
连续状态方程
\[
\dot{\bm x}(t)=A\bm x(t)+B u(t),\qquad y(t)=C\bm x(t),
\]
其中
\[
A=\begin{bmatrix}0&1\\0&0\end{bmatrix},\quad
B=\begin{bmatrix}0\\1\end{bmatrix},\quad
C=\begin{bmatrix}1&0\end{bmatrix}.
\]
采样周期为 $T$，零阶保持输入下
\[
\bm x(k+1)=\Phi\bm x(k)+\Gamma u(k),
\]
其中
\[
\Phi=e^{AT},\qquad
\Gamma=\int_0^T e^{A\tau}B\,d\tau.
\]
因为 $A^2=0$，
\[
e^{A\tau}=I+A\tau=
\begin{bmatrix}1&\tau\\0&1\end{bmatrix},
\]
故
\[
\Phi=
\begin{bmatrix}1&T\\0&1\end{bmatrix},
\]
并且
\[
\Gamma=
\int_0^T
\begin{bmatrix}1&\tau\\0&1\end{bmatrix}
\begin{bmatrix}0\\1\end{bmatrix}d\tau
=
\int_0^T\begin{bmatrix}\tau\\1\end{bmatrix}d\tau
=
\begin{bmatrix}T^2/2\\T\end{bmatrix}.
\]
因此离散状态方程为
\[
\boxed{
\bm x(k+1)=
\begin{bmatrix}1&T\\0&1\end{bmatrix}\bm x(k)
+
\begin{bmatrix}T^2/2\\T\end{bmatrix}u(k)
},
\qquad
\boxed{y(k)=\begin{bmatrix}1&0\end{bmatrix}\bm x(k)}.
\]

\section*{7-17}
已知
\[
\bm x(k+1)=A\bm x(k)+Bu(k),\qquad y(k)=C\bm x(k),
\]
其中
\[
A=\begin{bmatrix}1&1\\0&1\end{bmatrix},\quad
B=\begin{bmatrix}1/2\\1\end{bmatrix},\quad
C=\begin{bmatrix}1&0\end{bmatrix},\quad D=0.
\]
零初始条件下
\[
G(z)=C(zI-A)^{-1}B+D.
\]
有
\[
zI-A=
\begin{bmatrix}z-1&-1\\0&z-1\end{bmatrix},
\]
所以
\[
(zI-A)^{-1}=
\begin{bmatrix}
\dfrac{1}{z-1}&\dfrac{1}{(z-1)^2}\\[6pt]
0&\dfrac{1}{z-1}
\end{bmatrix}.
\]
于是
\[
\begin{aligned}
G(z)
&=
\begin{bmatrix}1&0\end{bmatrix}
\begin{bmatrix}
\dfrac{1}{z-1}&\dfrac{1}{(z-1)^2}\\
0&\dfrac{1}{z-1}
\end{bmatrix}
\begin{bmatrix}1/2\\1\end{bmatrix}\\
&=\frac{1}{2(z-1)}+\frac{1}{(z-1)^2}
=\frac{z+1}{2(z-1)^2}.
\end{aligned}
\]
因此
\[
\boxed{
G(z)=\frac{Y(z)}{U(z)}
=\frac{z+1}{2(z-1)^2}
=\frac{0.5z+0.5}{z^2-2z+1}
}.
\]
\end{document}
